\documentclass[10pt,a4paper]{amsart}
\usepackage[margin=19mm]{geometry}
\usepackage{amsmath,amssymb,mathtools}
\usepackage[scr=rsfs,cal=euler]{mathalfa}
\usepackage{xfrac}
\usepackage[unicode,hidelinks,bookmarks=false]{hyperref}
\newcommand{\ie}{i.e.\ }
\newcommand{\cf}{cf.\ }
\newcommand{\resp}{respectively\ }
\newcommand{\Subsection}{\subsection}
\providecommand{\nobreakdash}{\nobreak\hbox{-}}
\setlength{\emergencystretch}{2em}
\multlinegap0pt
\usepackage[shortlabels]{enumitem}
\usepackage{amssymb}
\usepackage{mathtools}
\newtheorem{theorem}{Theorem}
\DeclareMathOperator{\Cov}{Cov}
\newcommand{\E}{\mathbb E}
\newcommand{\N}{\mathbb N}
\DeclareMathOperator{\Var}{Var}
\DeclareMathOperator{\diag}{diag}
\newcommand{\one}{\mathbf 1}
\title{Growing Alphabets in Canonical Shuffle Experiments:\\ Likelihood-Ratio Laws, Estimation Bounds, and Low-Budget Equivariant Design}
\author{Alex Shvets}
\date{Source: arXiv:2603.18080v2 (CC BY 4.0)}
\begin{document}
\maketitle

\noindent\emph{Source and licence.} Growing Alphabets in Canonical Shuffle Experiments:\\ Likelihood-Ratio Laws, Estimation Bounds, and Low-Budget Equivariant Design. Version arXiv:2603.18080v2 (CC BY 4.0), distributed by arXiv under the Creative Commons Attribution 4.0 International licence, http://creativecommons.org/licenses/by/4.0/, as recorded in the arXiv licence metadata for this exact version and shown on the abstract page https://arxiv.org/abs/2603.18080v2. Full licence notice: \texttt{licenses/arxiv-2603.18080-CC-BY-4.0.txt}.\par\smallskip

\noindent\textit{Editorial scope.} Counted unit: the article's theorem thm:grr-mixtures (source0.tex lines 2297--2385). The article's complete original proof of it is delivered with it (source0.tex lines 2387--2568, one \texttt{proof} environment of 182 lines). No mathematics is written, repaired or re-derived: the statement and the proof are the article's own lines, in the article's own notation, and no retained line is substituted or re-flowed.  No further source-local context is needed: every cross-reference of every retained line resolves inside the two delivered spans.  Nothing was omitted from any retained span, so the package carries no omission record.  No retained line contains a citation, so the wrapper carries no bibliography and no bibliography entry.  No retained line contains a figure environment, an image-inclusion command or drawing code, so no image is delivered and the package declares no asset dependency.  Editorial formatting only, in addition: an AMS \texttt{amsart} preamble that restates the article's own declarations and macros used by the retained lines, namely the environments \texttt{theorem} on separate counters, together with the standard packages those lines need.  The retained environments print in this document as Theorem 1 in order of appearance, whereas the article prints the same items with the numbering of its own full document.  The complete native source of the article as distributed in the arXiv e-print for this exact version is bundled unchanged as \texttt{sources/196668/source0.tex}.
\begin{theorem}[GRR mixtures and null refinements]\label{thm:grr-mixtures}
Fix \(m\ge 1\).  Let
\[
\mathcal Y=\bigsqcup_{i=1}^m [d]_i \sqcup \mathcal Z,
\]
where \(\mathcal Z\) is an arbitrary finite set of null symbols.
For \(i=1,\dots,m\), choose \(p_i\ge 0\), \(\lambda_i>1\), and define
\[
W((i,y)\mid x)
=
p_i\left(
\beta_{\lambda_i}+\eta_{\lambda_i}\,\one\{y=x\}
\right),
\qquad
y\in[d]_i,
\]
while \(W(z\mid x)=r_z\) for \(z\in\mathcal Z\), with \(\sum_i p_i+\sum_{z\in\mathcal Z}r_z=1\).

Then the following hold.

\begin{enumerate}[label=\textup{(\roman*)}]
\item For every pair \(a\neq b\),
\[
\chi^2\!\bigl(W(\cdot\mid b)\,\|\,W(\cdot\mid a)\bigr)
=
\sum_{i=1}^m p_i C_{\lambda_i}. \tag{10.2}\label{eq:mix-chi}
\]

\item Assume \(S>0\) (equivalently, at least one block has \(p_i>0\)).
Let \(N^{(i)}\in\N^d\) be the block histogram on \([d]_i\), let
\[
M_i:=\sum_{y=1}^d N_y^{(i)},
\qquad
S:=\sum_{i=1}^m p_i\eta_{\lambda_i}^2,
\]
and define the affine-projected inverse estimator
\[
\widetilde\theta
:=
\frac{\one}{d}
+
\frac{1}{S}\sum_{i=1}^m \eta_{\lambda_i}
\left(
\frac{N^{(i)}}{n}-\frac{M_i}{nd}\,\one
\right). \tag{10.3}\label{eq:mix-estimator}
\]
Then \(\widetilde\theta\) is unbiased, \(\widetilde\theta-\theta\in T_d\), and its exact
fixed-composition risk is constant in \(\theta\):
\[
\E_\theta\|\widetilde\theta-\theta\|_2^2
=
\frac{d-1}{nd}\left(\frac{1}{S}-1\right). \tag{10.4}\label{eq:mix-risk}
\]

\item Among all channels of this form with prescribed canonical budget
\[
\sum_{i=1}^m p_i C_{\lambda_i}=C,
\]
the maximal signal coefficient \(S\) is
\[
S_{\mathrm{opt}}(C)
=
\begin{cases}
\displaystyle \frac{C}{d+2\sqrt{d-1}}, & 0\le C\le C_\ast(d),\\[1.2ex]
\displaystyle \eta_{\lambda(C)}^2, & C\ge C_\ast(d),
\end{cases} \tag{10.5}\label{eq:S-frontier}
\]
where
\[
\lambda(C)\ \text{ is the unique solution of }\ C_{\lambda(C)}=C,
\qquad
C_\ast(d):=C_{\sqrt{d-1}}.
\]
Consequently the optimal fixed-composition risk within this class is
\[
R^{\mathrm{fc}}_{\mathrm{opt}}(C)
=
\frac{d-1}{nd}\left(\frac{1}{S_{\mathrm{opt}}(C)}-1\right). \tag{10.6}\label{eq:mix-opt-risk}
\]
For \(0\le C\le C_\ast(d)\) it is attained by one-level augmented GRR with
\[
\lambda_\ast=\sqrt{d-1},
\qquad
p=\frac{C}{C_\ast(d)}. \tag{10.7}\label{eq:opt-low-budget}
\]
For \(C\ge C_\ast(d)\) it is attained by the calibrated ordinary GRR channel with \(p=1\) and
\(\lambda=\lambda(C)\).
\end{enumerate}
\end{theorem}

\begin{proof}
For \eqref{eq:mix-chi}, fix \(a\neq b\).
In block \(i\), the only nontrivial likelihood-ratio values are
\[
\lambda_i^{-1}\ \text{ on } (i,a),\qquad
\lambda_i\ \text{ on } (i,b),\qquad
1\ \text{ elsewhere},
\]
while every null symbol has ratio \(1\).  Summing the blockwise GRR contributions gives
\eqref{eq:mix-chi}.

For \eqref{eq:mix-risk}, fix \(u\in T_d\), and write
\[
\varphi_u((i,y)):=\frac{\eta_{\lambda_i}}{S}\,u_y,
\qquad
\varphi_u(z):=0,\quad z\in\mathcal Z.
\]
Because \(\sum_{y=1}^d u_y=0\),
\[
\langle u,\widetilde\theta\rangle
=
\frac{1}{n}\sum_{t=1}^n \varphi_u(Y_t).
\]
For an input \(x\),
\[
\E_x[\varphi_u(Y)]
=
\frac{1}{S}\sum_{i=1}^m p_i\eta_{\lambda_i}
\Bigl(\beta_{\lambda_i}\sum_{y=1}^d u_y+\eta_{\lambda_i}u_x\Bigr)
=
u_x
\]
because \(S=\sum_i p_i\eta_{\lambda_i}^2\).
Similarly,
\[
\E_x[\varphi_u(Y)^2]
=
\frac{1}{S^2}\sum_{i=1}^m p_i\eta_{\lambda_i}^2
\Bigl(\beta_{\lambda_i}\|u\|_2^2+\eta_{\lambda_i}u_x^2\Bigr).
\]
Hence
\[
\Var_x(\varphi_u(Y))
=
\frac{1}{S^2}\sum_{i=1}^m p_i\eta_{\lambda_i}^2
\Bigl(\beta_{\lambda_i}\|u\|_2^2+\eta_{\lambda_i}u_x^2\Bigr)-u_x^2.
\]
Averaging over the fixed composition \(\theta\) and writing
\(\Gamma_\theta:=\Cov_\theta(\widetilde\theta)\),
we obtain
\[
u^\top\Gamma_\theta u
=
\frac{1}{n}
\left[
\frac{\sum_i p_i\eta_{\lambda_i}^2\beta_{\lambda_i}}{S^2}\|u\|_2^2
+
\left(
\frac{\sum_i p_i\eta_{\lambda_i}^3}{S^2}-1
\right)
\sum_{x=1}^d \theta_x u_x^2
\right].
\]
Taking traces on \(T_d\) and using
\[
\operatorname{tr}_{T_d}(I)=d-1,
\qquad
\operatorname{tr}_{T_d}\!\bigl(P_T\diag(\theta)P_T\bigr)=\frac{d-1}{d},
\qquad
\beta_{\lambda_i}+\frac{\eta_{\lambda_i}}{d}=\frac1d,
\]
gives
\[
\operatorname{tr}_{T_d}\Gamma_\theta
=
\frac{d-1}{n}
\left[
\frac{1}{S^2}\sum_{i=1}^m p_i\eta_{\lambda_i}^2
\left(\beta_{\lambda_i}+\frac{\eta_{\lambda_i}}{d}\right)
-\frac1d
\right]
=
\frac{d-1}{nd}\left(\frac{1}{S}-1\right).
\]
Since \(\widetilde\theta-\theta\in T_d\), the total squared error equals its \(T_d\)-trace:
\[
\E_\theta\|\widetilde\theta-\theta\|_2^2
=
\operatorname{tr}_{T_d}\Gamma_\theta,
\]
which is exactly \eqref{eq:mix-risk}.

It remains to optimize \(S\) at fixed budget \(C\).
Define
\[
s(\lambda):=\eta_\lambda^2=\frac{(\lambda-1)^2}{(\lambda+d-1)^2},
\qquad
c(\lambda):=C_\lambda=\frac{(\lambda-1)^2(\lambda+1)}{\lambda(\lambda+d-1)}.
\]
Since
\[
c'(\lambda)
=
\frac{(\lambda-1)\bigl(2d\lambda^2+d\lambda+d+\lambda^3-\lambda^2+\lambda-1\bigr)}
{\lambda^2(\lambda+d-1)^2}
>0,
\]
the equation \(c(\lambda)=C\) has a unique solution \(\lambda(C)\) for every \(C>0\).

Every channel in the theorem yields a convex combination
\[
(C,S)=\sum_{i=1}^m p_i\bigl(c(\lambda_i),s(\lambda_i)\bigr)+(1-\sum_i p_i)(0,0).
\]
Hence the optimal frontier is the upper concave envelope of the planar curve
\(\lambda\mapsto (c(\lambda),s(\lambda))\) together with the origin.

First,
\[
\frac{s(\lambda)}{c(\lambda)}
=
\frac{\lambda}{(\lambda+1)(\lambda+d-1)}. \tag{10.8}\label{eq:ratio-sc}
\]
Its derivative is
\[
\frac{d}{d\lambda}\frac{s(\lambda)}{c(\lambda)}
=
\frac{d-1-\lambda^2}{(\lambda+1)^2(\lambda+d-1)^2}, \tag{10.9}\label{eq:ratio-derivative}
\]
so the slope from the origin is maximized uniquely at
\(
\lambda_\ast=\sqrt{d-1}
\).
Thus the line from the origin tangent to the curve meets it at
\[
\bigl(c(\lambda_\ast),s(\lambda_\ast)\bigr)
=
\left(
C_\ast(d),\frac{C_\ast(d)}{d+2\sqrt{d-1}}
\right).
\]

Second,
\[
\frac{ds}{dc}
=
\frac{2d\lambda^2}{(\lambda+d-1)\bigl(2d\lambda^2+d\lambda+d+\lambda^3-\lambda^2+\lambda-1\bigr)}, \tag{10.10}\label{eq:dsdc}
\]
and therefore
\[
\frac{d^2s}{dc^2}
=
\frac{2d\lambda^3 P_d(\lambda)}
{(\lambda-1)\bigl(2d\lambda^2+d\lambda+d+\lambda^3-\lambda^2+\lambda-1\bigr)^3}, \tag{10.11}\label{eq:d2sdc2}
\]
where
\[
P_d(\lambda)
=
d^2\lambda+2d^2-3d\lambda^3+d\lambda-4d-2\lambda^4+2\lambda^3-2\lambda+2.
\]
For \(\lambda\ge \sqrt{d-1}\),
\[
P_d'(\lambda)
=
d^2+d-2-(9d-6)\lambda^2-8\lambda^3
\le
-8(d-1)^2-8(d-1)^{3/2}<0, \tag{10.12}\label{eq:P-decreasing}
\]
and
\[
P_d(\sqrt{d-1})=2(d-1)^{3/2}(2-d)\le 0. \tag{10.13}\label{eq:P-at-star}
\]
Hence \(P_d(\lambda)\le 0\) for every \(\lambda\ge \sqrt{d-1}\), so
\[
\frac{d^2s}{dc^2}\le 0
\qquad\text{for } \lambda\ge \sqrt{d-1}. \tag{10.14}\label{eq:curve-concave}
\]
Thus the upper concave envelope consists exactly of the tangent line from the origin up to
\(C_\ast(d)\), followed by the original curve for \(C\ge C_\ast(d)\).
This proves \eqref{eq:S-frontier}, and \eqref{eq:mix-opt-risk} follows from \eqref{eq:mix-risk}.
The descriptions of the extremizers are immediate from the two pieces of the concave envelope.
\end{proof}

\end{document}
